\section{Deployable routers do not convert the signal}\label{sec:routing}

\paragraph{The comparison a router must win.} A router that sends some tasks to an expensive mode and others to a cheap one trades success for cost. Beating every fixed mode on both axes is too strict, but beating the fixed modes one at a time is too lenient: sending a random share of tasks to each of two modes is also a deployable policy. We therefore compare against the \emph{frontier of the fixed modes and all their random mixtures} --- the upper concave envelope of the fixed points in the (cost, success) plane, the non-decreasing convex hull that model-routing benchmarks use as a baseline \citep{hu2024routerbench,jitkrittum2025uniroute} --- and report how far above it a router reaches.

\paragraph{Routers.} Two constructions use 18 features available before a task runs: 14 intent-keyword indicators, the intent length, and three statistics of the first page. \emph{Per-arm heads} fit one logistic success model per mode and send a task to the cheapest mode whose predicted success clears a threshold. \emph{Triage} predicts whether any mode can solve the task, sends predicted-hopeless tasks to the cheapest mode and the rest to the best one. All predictions are out-of-fold (task-held-out five-fold).

\begin{table*}[t]
\centering\small
\setlength{\tabcolsep}{6pt}
\begin{tabular}{@{}llrr@{}}
\toprule
cells & policy & curve max (p) & deployable [5\%, 95\%] \\
\midrule
3 arms, 8 cells pooled & per-arm heads & +0.86 (0.072) & +0.23 [$-$2.20, +0.22] \\
 & triage & +1.69 (0.001) & +0.59 [$-$1.53, +1.23] \\
\addlinespace[2pt]
6 arms, 8 cells pooled & per-arm heads & +0.67 (0.252) & $-$0.38 [$-$3.20, $-$0.14] \\
 & triage & +1.24 (0.043) & $-$0.39 [$-$2.42, +0.81] \\
\addlinespace[2pt]
GPT-5.6, classifieds (5 arms) & per-arm heads & +0.45 (0.768) & $-$2.45 [$-$8.61, +4.24] \\
 & triage & +3.86 (0.127) & $-$1.63 [$-$6.70, +5.74] \\
\addlinespace[2pt]
Shopping B0, clean (3 arms) & per-arm heads & +1.14 (0.062) & $-$0.18 [$-$2.45, +2.02] \\
 & triage & +0.20 (0.496) & $-$0.87 [$-$2.63, +1.37] \\
\addlinespace[2pt]
Shopping B0, all tasks (3 arms) & per-arm heads & +1.09 (0.049) & +1.30 [$-$2.31, +1.85] \\
 & triage & +0.18 (0.541) & $-$0.05 [$-$2.14, +1.45] \\
\addlinespace[2pt]
Shopping B1, clean (6 arms) & per-arm heads & +1.11 (0.143) & +0.56 [$-$4.69, +0.88] \\
 & triage & +1.63 (0.158) & +1.19 [$-$4.10, +1.57] \\
\addlinespace[2pt]
Shopping B1, all tasks (6 arms) & per-arm heads & +0.48 (0.169) & $-$0.43 [$-$2.78, +0.73] \\
 & triage & +0.69 (0.289) & +0.40 [$-$2.02, +2.02] \\
\bottomrule
\end{tabular}
\caption{Routing gain over the fixed arms and their random mixtures, in SR points. \emph{Curve max}: the best point of an out-of-fold curve, chosen after seeing test outcomes, with its label-shuffle $p$ (pooled rows: max over a normalised budget). \emph{Deployable}: the operating point chosen on training folds only, task bootstrap. Sources: \texttt{routing\_three\_arm}, \texttt{representation\_routing\_frontier}, \texttt{routing\_gain\_upper\_bounds}, \texttt{routing\_extension\_cells}.}
\label{tab:routing}
\end{table*}


\paragraph{Room above the frontier exists, and it is difficulty.} Sweeping each router's threshold traces a curve. Over the three deployment modes, the triage curve pooled over the eight cells rises 1.69 points above the frontier at low budget, which a label-shuffle null that refits the whole curve does not produce ($p=0.001$; Table~\ref{tab:routing}). The per-arm heads reach 0.86 points ($p=0.072$). Triage does not choose a representation for a task; it decides which tasks are not worth the expensive mode. The room that exists is the difficulty lever, not representation fit.

\paragraph{A deployable router does not capture it.} A curve's best point is chosen after seeing test outcomes. When the operating point is instead chosen on training folds only, the pooled gain is +0.23 points for per-arm heads and +0.59 for triage, and neither interval excludes zero (Table~\ref{tab:routing}). Over all six modes the deployable per-arm router sits significantly \emph{below} the frontier ($-0.38$; 95\% upper bound $-0.14$). % routing_gain_upper_bounds pooled
The pattern holds for the strongest backbone: for GPT-5.6 on classifieds, where the screenshot's value was largest (\S\ref{sec:screenshot}), the deployable routers lose 1.6--2.5 points; on both shopping cells they stay within 1.2 points of the frontier.

\paragraph{Robustness.} Holding tasks of the same template out together lowers the six-mode per-arm router's mean deployable gain from $-0.38$ to $-1.60$ points, and fitting on one run and scoring on its rerun gives point estimates of both signs whose intervals all include zero (Appendix~\ref{app:robust}). % routing_crossrun_template_validation
Repricing cost as GPU time or wall-clock time leaves the frontier result unchanged, but changes \emph{which} single cell's triage curve passes a per-cell Holm correction, so no single-cell result is quoted here. % §545
One of the three-mode intervals (per-arm heads, pooled) has an upper end below its point estimate, a sign that the bootstrap of a policy with a selected threshold is biased; we read these intervals as magnitudes, not as tests.
